\section{Background: Composable Timing-Sensitive Non-Interference}\label{sec:background}

This section recalls the calculus and definitions of Rafnsson et
al.~\cite{rafnsson2017}, in the form we mechanise. The form differs from the
original in four places, discussed in \cref{sec:departures}.

\subsection{Processes, traces and time}

A process $p : \Proc{I}{O}$ has an input type $I$ and an output type $O$. It moves
in one of two ways, by receiving an input ($p \xrightarrow{?i} p'$) or by emitting
an output ($p \xrightarrow{!o} p'$), in any order. Both kinds of step are total and
deterministic: in every state a process can receive any input and can emit, and the
successor (and on an output step the value emitted) is determined. A \emph{trace} is
the sequence of labels along a finite run.

Time is counted in output steps. Every output step is one tick of a clock the
observer shares, and on every tick a process emits a value; there is no silent tick.
Where nothing happens, the value says so (in our models a component that did not run
emits \code{None}), and that absence is itself observable. Input steps take no time. This
is what makes the setting timing-sensitive: delaying an output by one step changes
the trace exactly as changing its value would.

Processes are built from seven constructors (\cref{tab:constructors}). The
operating system of \cref{sec:model} is a single closed term over them.

\begin{table}[t]
\centering
\caption{The constructors, by their behaviour on an input and on an output step.}
\label{tab:constructors}
\small
\begin{tabularx}{\columnwidth}{@{}lL@{}}
\toprule
constructor & behaviour \\
\midrule
\code{out o} & receives anything and is unchanged; emits the constant \code{o} \\
\code{map f g p} & hands \code{p} the input \code{f i}; emits \code{g o} when \code{p} emits \code{o} \\
\code{sta f g v p} & a state cell \code{v}: on input \code{i} the cell becomes \code{f i v} and \code{p}
  receives the pair; on output \code{o} of \code{p} the cell becomes \code{g o v} and the pair
  of new cell and \code{o} is emitted \\
\code{swi b p} & a one-bit switch: in phase \code{false} it emits \code{None} and \code{p} does not
  step; in phase \code{true} it emits \code{Some o} for \code{p}'s tagged output \code{(c,o)}, and
  the tag \code{c} and the inputs flip the phase \\
\code{par p1 p2} & both receive every input; both emit, and the pair is emitted \\
\code{loop p} & feedback: \code{p} emits \code{o}, \code{o} is fed back to \code{p} as an input, and \code{o}
  is emitted \\
\code{maybe p} & optional input: \code{p} ignores \code{None} and receives \code{i} on \code{Some i} \\
\bottomrule
\end{tabularx}
\end{table}

\subsection{Characterised equivalences and non-interference}

Security is stated relative to a lattice of levels. Everything in this paper can be
read with two: the attacker $\bot$, and $\top$, who sees everything
(\cref{sec:generality} discusses arbitrary lattices). Each input and output type
carries a \emph{characterised equivalence} (\m{cEquiv}), a record of two
level-indexed predicates:
\begin{itemize}
\item $\rel\;l\;x\;y$: an observer at $l$ cannot tell $x$ from $y$;
\item $\dis\;l\;x$: an observer at $l$ does not see $x$ at all, so $x$ may be
  inserted or removed freely.
\end{itemize}
$\rel\;l$ is an equivalence, and $\dis\;l$ is \emph{characterised} as one of its
classes:
\[
\dis\;l\;a_0 \;\Rightarrow\; \forall a_1.\;
  (\dis\;l\;a_1 \Leftrightarrow \rel\;l\;a_0\;a_1).
\]
So at each level the type is partitioned into classes the observer cannot tell
apart, and one of them, the \emph{distinguished class}, holds the values the
observer does not see. It is empty when the level hides nothing.

Let $\mathit{insert}\;n\;i\;t$ place input $i$ at position $n$ of trace $t$, and let
$\mathsf{Trace}$ compare outputs up to $\rel_{O}\;l$.

\begin{definition}[Non-interference]
$\NI\;\mathit{IRel}\;\mathit{ORel}\;p$ holds when for every level $l$, every trace
$t$, position $n$ and inputs $i, i'$:
\begin{enumerate}
\item (substitution) if $\rel_I\;l\;i\;i'$ and $\mathit{insert}\;n\;i\;t$ is a trace
  of $p$, then so is $\mathit{insert}\;n\;i'\;t$;
\item (insertion) if $\dis_I\;l\;i$ and $t$ is a trace of $p$, then so is
  $\mathit{insert}\;n\;i\;t$;
\item (deletion) if $\dis_I\;l\;i$ and $\mathit{insert}\;n\;i\;t$ is a trace of
  $p$, then so is $t$.
\end{enumerate}
\end{definition}

The set of traces visible at $l$ is thus unchanged by anything the observer at $l$
may not see. All the content is at $\bot$: at $\top$ nothing is unobservable and
indistinguishable means equal, so the clauses hold trivially.

\subsection{Equivalences}

Two equivalences do most of the work. Under \m{public_equiv} nothing is unobservable
and values are related iff equal, at every level. Under \m{private_equiv} every value
is unobservable at $\bot$, where all values are related, and the equivalence is
\m{public_equiv} at $\top$.

Composite equivalences are built from these. \m{eqpair} relates pairs componentwise,
and its variants \m{eqpair_L}, \m{eqpair_R} and \m{eqpair_LR} make the pair secret when
its left, right, or both components are. \m{eqsum} relates a sum tag by tag and never
relates an \code{inl} to an \code{inr}, so the tag stays public even when the payload is not;
\m{eqsum_L} makes an \code{inl} secret exactly when its payload is. The \m{eqmaybe} family
relates \code{Option} values, deferring to an inner equivalence on \code{Some} and differing
in who can see \code{None} (\cref{tab:eqmaybe}). By the characterisation law, \code{Some v}
and \code{None} are related at $l$ exactly when both are unobservable there.

\begin{table}[t]
\centering
\caption{Who can see \texttt{None} under the \texttt{eqmaybe} family. To the
  attacker, \texttt{eqmaybe\_private} and \texttt{eqmaybe\_hidden} agree.}
\label{tab:eqmaybe}
\small
\begin{tabularx}{\columnwidth}{@{}lLL@{}}
\toprule
equivalence & $\bot$ sees \texttt{None} & $\top$ sees \texttt{None} \\
\midrule
\m{eqmaybe}         & yes & yes \\
\m{eqmaybe_private} & no  & yes \\
\m{eqmaybe_hidden}  & no  & no \\
\code{eqmaybe_swi BRel} & iff $\bot$ is aware under \code{BRel} & iff $\top$ is aware \\
\bottomrule
\end{tabularx}
\end{table}

\subsection{Compositional theorems}\label{sec:composition}

The calculus comes with one non-interference theorem per constructor
(\cref{tab:composition}). Each derives the equivalences of the composite from those
of its parts, which is what gives the equivalences of a large model their shape.
Two carry substantial side conditions. \m{sta_NI} asks that the state update map
related states to related states,
\[
\mathit{fv\_NI}:\quad \rel\;l\;i\;i' \wedge \rel\;l\;v\;v' \;\Rightarrow\;
  \rel\;l\;(f\;i\;v)\;(f\;i'\;v'),
\]
and that an unobservable input not visibly move the state,
$\mathit{f\_EP}:\; \dis\;l\;i \Rightarrow \rel\;l\;(f\;i\;v)\;v$. \m{swi_NI} asks, at
each level, that the observer either can trust the gate bit (it is \emph{aware}) or
learns nothing from the gated process because all its outputs fall in one class (it
is \emph{oblivious}).

\begin{table*}[t]
\centering
\caption{The composition theorems, slightly simplified. $I$ and $O$ are the input
  and output equivalences, $V$ the state equivalence and $B$ the gate equivalence.}
\label{tab:composition}
\footnotesize
\begin{tabularx}{\textwidth}{@{}lLL@{}}
\toprule
theorem & from & to \\
\midrule
\m{out_NI} & --- & $\NI\;I\;O\;(\mathsf{out}\;o)$ \\
\m{map_NI} & $\NI\;I'\;O\;p$, side conditions on $f$ and $g$
  & $\NI\;I\;O'\;(\mathsf{map}\;f\;g\;p)$ \\
\m{sta_NI} & $\NI\;(\mathsf{eqpair\_R}\;V\;I)\;O\;p$, $\mathit{fv\_NI}$,
  $\mathit{f\_EP}$
  & $\NI\;I\;(\mathsf{eqpair}\;V\;O)\;(\mathsf{sta}\;f\;g\;v\;p)$ \\
\m{swi_NI} & $\NI\;I\;(\mathsf{eqpair\_LR}\;B\;O)\;p$, aware or oblivious at
  each level
  & $\NI\;(\mathsf{eqpair\_LR}\;B\;I)$\newline
    \hspace*{1em}$(\mathsf{eqmaybe\_swi}\;O\;B)\;(\mathsf{swi}\;b\;p)$ \\
\m{par_NI} & $\NI\;I\;O_1\;p_1$ and $\NI\;I\;O_2\;p_2$
  & $\NI\;I\;(\mathsf{eqpair}\;O_1\;O_2)\;(\mathsf{par}\;p_1\;p_2)$ \\
\m{loop_NI} & $\NI\;I\;I\;p$ & $\NI\;I\;I\;(\mathsf{loop}\;p)$ \\
\m{maybe_NI} & $\NI\;I\;O\;p$
  & $\NI\;(\mathsf{eqmaybe\_hidden}\;I)\;O\;(\mathsf{maybe}\;p)$ \\
\bottomrule
\end{tabularx}
\end{table*}
